% Hey, if you're using this preamble it means that it was probably written by Stefano Graziosi (me). If you see something that doesn't make sense, feel free to email me at stefano.graziosi@studbocconi.it
% p.s. in case it's not already evident from the preamble, I'm not a professional LaTeX user, so I'm sure there are better ways to do things. I'm just trying to make it work.

% =============================================================================
%  LAST UPDATE: 16-08-2026
% =============================================================================

% I don't own copyright on anything; credits to original authors.

% =============================================================================
%  1. ENCODING, FONTS & PAGE LAYOUT
% =============================================================================

% Encoding and fonts
\usepackage[T1]{fontenc}
\usepackage{lmodern}
\usepackage{mathrsfs}

% Page layout and document formatting
\usepackage[a4paper,margin=1in]{geometry}
\usepackage{fancyhdr}
\usepackage{setspace}
\usepackage{multicol}
\usepackage{enumitem}
\usepackage[normalem]{ulem}

% Titles and PDF inclusion
\usepackage{titlesec}
\usepackage{titlepic}
\usepackage[table,xcdraw,dvipsnames]{xcolor}
\usepackage{pdfpages}

% Dummy text / extras
\usepackage{kantlipsum}

% =============================================================================
%  2. COLORS & BOXES
% =============================================================================

\usepackage[many]{tcolorbox}

% Custom palette
\definecolor{sgblue}{RGB}{0,169,211}
\definecolor{sggreen}{RGB}{0,164,0}
\definecolor{sgpurple}{RGB}{99,0,165}
\definecolor{sgyellow}{RGB}{255,211,0}
\definecolor{sgorange}{RGB}{255,127,20}
\definecolor{sbblue}{RGB}{219,248,254}
\definecolor{sbgreen}{RGB}{223,255,218}
\definecolor{sbpurple}{RGB}{241,220,255}

% =============================================================================
%  3. MATHEMATICS, SYMBOLS & THEOREM INFRASTRUCTURE
% =============================================================================

\usepackage{amsmath}
\usepackage{amssymb}
\usepackage{amsthm}
\usepackage{mathtools}
\usepackage{amsfonts}
\usepackage{bm}
\usepackage{thmtools}
\usepackage{stmaryrd}
\usepackage[version=4]{mhchem}
\usepackage{halloweenmath}

% =============================================================================
%  4. TABLES
% =============================================================================

\usepackage{array}
\usepackage{tabularx}
\usepackage{booktabs}
\usepackage{threeparttable}
\usepackage{multirow}
\usepackage{siunitx}

\newcolumntype{Y}{>{\raggedright\arraybackslash}X}

\sisetup{
    round-mode=places,
    round-precision=3,
    table-number-alignment=center
}

% =============================================================================
%  5. FIGURES, DIAGRAMS & PLOTS
% =============================================================================

\usepackage{graphicx}
\usepackage{wrapfig}
\usepackage{float}
\usepackage{caption}
\usepackage{subcaption}
\usepackage[export]{adjustbox}

\captionsetup{
    font=small,
    labelfont=bf,
    labelsep=period
}

% Diagrams
\usepackage[all]{xy}
\usepackage{tikz}
\usetikzlibrary{calc,arrows.meta,positioning}

% Optional speed-up for heavy documents; requires -shell-escape
% \usetikzlibrary{external}
% \tikzexternalize[prefix=tikzcache/]

% Plots
\usepackage{pgfplots}
\pgfplotsset{compat=1.18}
\usepgfplotslibrary{groupplots}

% =============================================================================
%  6. VERBATIM & CODE TOOLS
% =============================================================================

\usepackage{listings}
\usepackage{fvextra}
\usepackage{codehigh}

% =============================================================================
%  7. QUOTATIONS
% =============================================================================

\usepackage{csquotes}

% =============================================================================
%  8. HYPERLINKS, BOOKMARKS & CROSS-REFERENCES
% =============================================================================

% Keep hyperlink-related packages late in the package-loading sequence.
\usepackage{hyperref}
\hypersetup{
    pdftitle={\jobname},
    pdfauthor={Stefano Graziosi},
    hypertexnames=false,
    colorlinks=true,
    linkcolor=MidnightBlue,
    citecolor=ForestGreen,
    urlcolor=sgpurple
}

\usepackage{bookmark}
\usepackage[nameinlink,noabbrev]{cleveref}

% =============================================================================
%  9. NUMBERING
% =============================================================================

% Equations numbered within sections
\numberwithin{equation}{section} % or \counterwithin{equation}{section}

% Figures and tables numbered within sections
\usepackage{chngcntr}
\counterwithin{figure}{section}
\counterwithin{table}{section}
\renewcommand{\thefigure}{\thesection-\arabic{figure}}
\renewcommand{\thetable}{\thesection-\arabic{table}}

% Keep running heads compact even when chapter and section titles are long. 
\renewcommand{\footrulewidth}{0pt}  
% Article-like classes do not define \chaptermark, so only customize it when
% the selected document class supports chapters.
\makeatletter
\@ifundefined{chaptermark}{}{%
    \renewcommand{\chaptermark}[1]{\markboth{Chapter~\thechapter}{}}%
}
\@ifundefined{sectionmark}{%
    \renewcommand{\sectionmark}[1]{\markright{\thesection\quad #1}}
}{%
    \renewcommand{\sectionmark}[1]{\markright{\thesection\quad #1}}
}
\makeatother
\setlength{\headheight}{14pt}

% Give dense prose a little extra flexibility and avoid stretched page bottoms.
\setlength{\emergencystretch}{3em}
\raggedbottom

% =============================================================================
%  10. GENERIC BOX ENVIRONMENT
% =============================================================================

% Standard LaTeX box with color accents
\newtcolorbox{mybox}[3][]{
    colframe=#2!25,
    colback=#2!10,
    coltitle=#2!20!black,
    title={#3},
    #1,
}

% =============================================================================
%  11. BASIC PROBLEM / SOLUTION ENVIRONMENTS
% =============================================================================

% "Problem" environment
\newtheorem{problem}{Problem}

% Personalized "Solution" environment
\newenvironment{solution}[1][\it{\textcolor{MidnightBlue}{Solution}}]
    {\textbf{#1. } }
    {\textcolor{MidnightBlue}{$\square$}}

% =============================================================================
%  12. THEOREM STYLES & BOXED PRESENTATION
% =============================================================================

% Helper length for theorem-heading alignment.
\newlength{\spacelength}
\settowidth{\spacelength}{\normalfont\ }

% Common tcolorbox settings. These only remove repeated code; the individual
% environments below retain the same colours, borders, backgrounds and spacing.
\tcbset{
    sg theorem base/.style={
        boxrule=0pt,
        boxsep=0pt,
        before skip=10pt,
        after skip=10pt,
        left=5pt,
        right=5pt,
        breakable,
    },
    sg theorem filled/.style={
        sg theorem base,
        enhanced jigsaw,
        sharp corners,
    },
    sg theorem blank/.style={
        blanker,
        sg theorem base,
        sharp corners,
    },
}

% -----------------------------------------------------------------------------
% Main theorem style
% -----------------------------------------------------------------------------

\declaretheoremstyle[
    headfont={\bfseries\sffamily\footnotesize},
    notefont={\normalfont},
    bodyfont={\normalfont},
    headpunct={\relax},
    headformat={%
        \makebox[0pt][r]{\NAME\ \NUMBER\hspace{\marginparsep}}%
        \hskip-\spacelength{\normalsize\NOTE}%
    },
]{theorem}

% thmtools' \declaretheoremstyle is a front-end to the theorem-style backend
% supplied by amsthm, so this style can also be selected with \theoremstyle.

% Theorem
\declaretheorem[style=theorem]{theorem}
\tcolorboxenvironment{theorem}{
    sg theorem filled,
    colback={White!94!ForestGreen},
    borderline west={1pt}{0pt}{ForestGreen},
}

% Proposition
\declaretheorem[style=theorem]{proposition}
\tcolorboxenvironment{proposition}{
    sg theorem filled,
    colback={White!94!Mulberry},
    borderline west={1pt}{0pt}{Mulberry},
}

% -----------------------------------------------------------------------------
% Proof style
% -----------------------------------------------------------------------------

% amsthm already defines proof, so neutralize it before redeclaring it.
\let\proof\relax
\let\endproof\relax

\declaretheoremstyle[
    headfont={\small\scshape},
    notefont={\normalfont},
    bodyfont={\normalfont},
    headpunct={\relax},
    headformat={%
        \makebox[0pt][r]{\NAME\hspace{\marginparsep}}%
        \hskip-\spacelength{\NOTE}%
    },
]{proof}

\tcolorboxenvironment{proof}{
    blanker,
    sg theorem base,
    borderline west={1pt}{0pt}{black},
}

\declaretheorem[style=proof,qed=\qedsymbol]{proof}

% -----------------------------------------------------------------------------
% Claim / intuition style
% -----------------------------------------------------------------------------

\declaretheoremstyle[
    headfont={\footnotesize\itshape},
    notefont={\normalfont},
    bodyfont={\normalfont},
    headpunct={\relax},
    headformat={%
        \makebox[0pt][r]{\NAME\hspace{\marginparsep}}%
        \hskip-\spacelength{\NOTE}%
    },
]{claim}

% Capital-I Intuition environment (kept for backwards compatibility)
\declaretheorem[style=claim]{Intuition}

% -----------------------------------------------------------------------------
% Additional theorem-like environments
% -----------------------------------------------------------------------------

% Since the styles above are registered with the amsthm backend, selecting the
% "theorem" style here is intentional and valid.
\theoremstyle{theorem}

% Numbered Question (unboxed)
\newtheorem{ques}{Question}

% Definition
\newtheorem{definition}{Definition}
\tcolorboxenvironment{definition}{
    sg theorem filled,
    colback={White!94!Cerulean},
    borderline west={1pt}{0pt}{Cerulean},
}

% Lemma
\newtheorem{lemma}{Lemma}
\tcolorboxenvironment{lemma}{
    sg theorem blank,
    borderline west={1pt}{0pt}{Rhodamine},
}

% Remark
\newtheorem{remark}{Remark}
\tcolorboxenvironment{remark}{
    sg theorem filled,
    colback={White!94!BurntOrange},
    borderline west={1pt}{0pt}{BurntOrange},
}

% Hint
\newtheorem{hint}{Hint}
\tcolorboxenvironment{hint}{
    sg theorem filled,
    colback={White},
    borderline west={1pt}{0pt}{BurntOrange},
}

% Warning
\newtheorem{warning}{Warning}
\tcolorboxenvironment{warning}{
    sg theorem filled,
    colback={White},
    borderline west={1pt}{0pt}{red},
}

% Disclaimer
\newtheorem{disclaimer}{Disclaimer}
\tcolorboxenvironment{disclaimer}{
    sg theorem filled,
    colback={White},
    borderline west={1pt}{0pt}{red},
}

% Unnumbered Question
\newtheorem*{question}{Question}
\tcolorboxenvironment{question}{
    sg theorem filled,
    colback={White},
    borderline west={1pt}{0pt}{red},
}

% Corollary
\newtheorem{corollary}{Corollary}
\tcolorboxenvironment{corollary}{
    sg theorem filled,
    borderline west={1pt}{0pt}{WildStrawberry},
}

% Example
\newtheorem{example}{Example}
\tcolorboxenvironment{example}{
    sg theorem blank,
    borderline west={1pt}{0pt}{Dandelion},
}

% Lowercase helper environments (kept for backwards compatibility)
\theoremstyle{claim}
\newtheorem{intu}{Intuition}
\newtheorem{solu}{Solution}

% =============================================================================
%  CODE LISTINGS - GENERAL
% =============================================================================

% 1. Define your custom colors (Ensuring they are defined)
\definecolor{codegreen}{rgb}{0,0.6,0}
\definecolor{codegray}{rgb}{0.5,0.5,0.5}
\definecolor{codepurple}{rgb}{0.58,0,0.82}
\definecolor{backcolour}{rgb}{0.95,0.95,0.92}
% A specific color for Stata functions (to differentiate from commands)
\definecolor{statafunction}{rgb}{0.2,0.2,0.7} 

% 2. The Improved Stata Language Definition
\lstdefinelanguage{Stata}{
    sensitive=true,
    morecomment=[l]{//},
    morecomment=[s]{/*}{*/},
    morestring=[b]",
    morestring=[d]",
    morestring=[b]{`}{'},
    % Group 1: System/Flow (if, else, foreach, program, python)
    morekeywords=[1]{
        if, else, in, using, of, varlist, everything, foreach, forvalues, 
        while, continue, break, program, syntax, return, ereturn, sreturn, 
        args, capture, noisily, quietly, set, version, global, local, scalar, 
        matrix, mat, frame, python, java, putexcel, collect, do, run, ado, 
        clear, exit,
    },
    % Group 2: Data Manipulation (gen, egen, sort, merge)
    morekeywords=[2]{
        generate, gen, replace, egen, rename, ren, drop, keep, sort, bysort, 
        order, collapse, merge, append, reshape, encode, decode, destring, 
        tostring, recode, mvencode, mvdecode, use, save, import, export, 
        infile, infix, describe, desc, list, count, inspect, codebook, label, 
        notes, compress, format, cast,
    },
    % Group 3: Statistics (reg, sum, mixed, lasso, bayes)
    morekeywords=[3]{
        regress, reg, summarize, sum, tabulate, tab, logit, probit, tobit, 
        poisson, nbreg, xtreg, xtlogit, xtpoisson, xtset, xtdescribe, mixed, 
        meqrlogit, mepoisson, ivregress, areg, rreg, qreg, lasso, elasticnet, 
        bayes, bayesmh, meta, predict, predictnl, margins, marginsplot, test, 
        testnl, lincom, nlcom, contrast, pwcompare, estat, estimates, est, 
        pwcorr, correlate, corr, ttest, anova, oneway,
    },
    % Group 4: Graphics (twoway, scatter)
    morekeywords=[4]{
        graph, twoway, scatter, line, connected, bar, hbar, box, hbox, pie, 
        histogram, kdensity, function, title, subtitle, xtitle, ytitle, 
        legend, xlabel, ylabel, note, caption, saving,
    },
    % Group 5: Reserved/Types (byte, double, _n)
    morekeywords=[5]{
        byte, int, long, float, double, str, strL, c, _N, _n, _rc, _all, 
        aggregate, array, boolean, class, colvector, complex, const, delegate, 
        delete, eltypedef, enum, explicit, external, friend, inline, mata, 
        namespace, new, numeric, NULL, operator, orgtypedef, pointer, 
        polymorphic, pragma, private, protected, public, quad, real, 
        rowvector, short, signed, static, struct, super, switch, template, 
        this, throw, transmorphic, try, typedef, typename, union, unsigned, 
        vector, virtual, void, volatile,
    },
    % Groups 6-12: Functions (Dates, Math, Matrix, Strings)
    morekeywords=[6]{
        bofd, Cdhms, Chms, Clock, clock, Cmdyhms, Cofc, cofC, Cofd, cofd, 
        daily, date, day, dhms, dofb, dofC, dofc, dofh, dofm, dofq, dofw, 
        dofy, dow, doy, halfyear, halfyearly, hh, hhC, hms, hofd, hours, 
        mdy, mdyhms, minutes, mm, mmC, mofd, month, monthly, msofhours, 
        msofminutes, msofseconds, qofd, quarter, quarterly, seconds, ss, 
        ssC, tC, tc, td, th, tm, tq, tw, week, weekly, wofd, year, yearly, 
        yh, ym, yofd, yq, yw,
    },
    morekeywords=[7]{
        abs, ceil, cloglog, comb, digamma, exp, expm1, floor, invcloglog, 
        invlogit, ln, ln1m, ln1p, lnfactorial, lngamma, log, log10, log1m, 
        log1p, logit, max, min, mod, reldif, round, sign, sqrt, trigamma, 
        trunc, rbeta, rbinomial, rcauchy, rchi2, rexponential, rgamma, 
        rhypergeometric, rigaussian, rlaplace, rlogistic, rnbinomial, 
        rnormal, rpoisson, rt, runiform, runiformint, rweibull, rweibullph, 
        tin, twithin,
    },
    morekeywords=[8]{
        cholesky, coleqnumb, colnfreeparms, colnumb, colsof, det, diag, 
        diag0cnt, el, get, hadamard, I, inv, invsym, issymmetric, J, 
        matmissing, matuniform, mreldif, nullmat, roweqnumb, rownfreeparms, 
        rownumb, rowsof, sweep, trace, vec, vecdiag,
    },
    morekeywords=[9]{
        abbrev, char, collatorlocale, collatorversion, indexnot, plural, 
        regexm, regexr, regexs, soundex, soundex_nara, strcat, strdup, 
        string, stritrim, strlen, strlower, strltrim, strmatch, strofreal, 
        strpos, strproper, strreverse, strrpos, strrtrim, strtoname, strtrim, 
        strupper, subinstr, subinword, substr, tobytes, uchar, udstrlen, 
        udsubstr, uisdigit, uisletter, ustrcompare, ustrcompareex, ustrfix, 
        ustrfrom, ustrinvalidcnt, ustrleft, ustrlen, ustrlower, ustrltrim, 
        ustrnormalize, ustrpos, ustrregexm, ustrregexra, ustrregexrf, 
        ustrregexs, ustrreverse, ustrright, ustrrpos, ustrrtrim, 
        ustrsortkey, ustrsortkeyex, ustrtitle, ustrto, ustrtohex, 
        ustrtoname, ustrtrim, ustrunescape, ustrupper, ustrword, 
        ustrwordcount, usubinstr, usubstr, word, wordbreaklocale, wordcount,
    },
    morekeywords=[10]{
        acos, acosh, asin, asinh, atan, atanh, cos, cosh, sin, sinh, tan, 
        tanh,
    },
    morekeywords=[11]{
        betaden, binomial, binomialp, binomialtail, binormal, cauchy, 
        cauchyden, cauchytail, chi2, chi2den, chi2tail, dunnettprob, 
        exponential, exponentialden, exponentialtail, F, Fden, Ftail, 
        gammaden, gammap, gammaptail, hypergeometric, hypergeometricp, 
        ibeta, ibetatail, igaussian, igaussianden, igaussiantail, 
        invbinomial, invbinomialtail, invcauchy, invcauchytail, invchi2, 
        invchi2tail, invdunnettprob, invexponential, invexponentialtail, 
        invF, invFtail, invgammap, invgammaptail, invibeta, invibetatail, 
        invigaussian, invigaussiantail, invlaplace, invlaplacetail, 
        invlogistic, invlogistictail, invnbinomial, invnbinomialtail, 
        invnchi2, invnF, invnFtail, invnibeta, invnormal, invnt, invnttail, 
        invpoisson, invpoissontail, invt, invttail, invtukeyprob, 
        invweibull, invweibullph, invweibullphtail, invweibulltail, laplace, 
        laplaceden, laplacetail, lncauchyden, lnigammaden, lnigaussianden, 
        lniwishartden, lnlaplaceden, lnmvnormalden, lnnormal, lnnormalden, 
        lnwishartden, logistic, logisticden, logistictail, nbetaden, 
        nbinomial, nbinomialp, nbinomialtail, nchi2, nchi2den, nchi2tail, 
        nF, nFden, nFtail, nibeta, normal, normalden, npnchi2, npnF, npnt, 
        nt, ntden, nttail, poisson, poissonp, poissontail, t, tden, ttail, 
        tukeyprob, weibull, weibullden, weibullph, weibullphden, 
        weibullphtail, weibulltail,
    },
    morekeywords=[12]{
        autocode, byteorder, _caller, chop, clip, cond, e, fileexists, 
        fileread, filereaderror, filewrite, fmtwidth, has_eprop, inlist, 
        inrange, irecode, maxbyte, maxdouble, maxfloat, maxint, maxlong, mi, 
        minbyte, mindouble, minfloat, minint, minlong, missing, r, recode, 
        replay, smallestdouble,
    }
}

% 3. Your Style (Updated to include the new groups)
\lstdefinestyle{mystyle}{
    backgroundcolor=\color{backcolour},
    commentstyle=\color{codegreen},
    stringstyle=\color{codepurple},
    numberstyle=\tiny\color{codegray},
    basicstyle=\ttfamily\footnotesize,
    breakatwhitespace=false,
    breaklines=true,
    captionpos=b,
    keepspaces=true,
    numbers=left,
    numbersep=5pt,
    showspaces=false,
    showstringspaces=false,
    showtabs=false,
    tabsize=2,
    % --- KEYWORD MAPPING ---
    % Groups 1-4 are COMMANDS (regress, gen, etc.) -> Use your Magenta
    keywordstyle=[1]\color{magenta},
    keywordstyle=[2]\color{magenta},
    keywordstyle=[3]\color{magenta},
    keywordstyle=[4]\color{magenta},
    % Group 5 are TYPES (float, int) -> Use Blue or Magenta (User choice)
    keywordstyle=[5]\color{blue}, 
    % Groups 6-12 are FUNCTIONS (sum(), max()) -> Use a distinct color 
    % (Using Magenta here looks confusing as they look like commands). 
    % I selected 'statafunction' (defined above) which is a nice blue/purple.
    keywordstyle=[6]\color{statafunction},
    keywordstyle=[7]\color{statafunction},
    keywordstyle=[8]\color{statafunction},
    keywordstyle=[9]\color{statafunction},
    keywordstyle=[10]\color{statafunction},
    keywordstyle=[11]\color{statafunction},
    keywordstyle=[12]\color{statafunction},
}

\lstset{style=mystyle}